\documentclass[10pt,a4paper]{article}
\usepackage{xeCJK}
\usepackage{fontspec}
\usepackage{booktabs}
\usepackage{array}
\usepackage{xcolor}
\usepackage{geometry}
\usepackage[protrusion=true]{microtype}
\geometry{margin=1.7cm}
\linespread{1.05}

\setmainfont{Baskerville}
\setCJKmainfont{Songti SC}
\newfontfamily{\starfont}{Songti SC}

\definecolor{wrongred}{RGB}{198,40,40}
\definecolor{rulegray}{RGB}{190,190,190}

\newcommand{\checkmarkbox}{\fbox{\rule{0pt}{0.75ex}\rule{0.75ex}{0pt}}}
\newcommand{\STAR}{\textcolor{wrongred}{\starfont ★}}
\newcommand{\STARSLOT}[1]{\makebox[1.5em][l]{#1}}
% \qa{题号}{星标}{答案}{解析}
\newcommand{\qa}[4]{\noindent\STARSLOT{#2}\textbf{#1.}~\textbf{#3}\par\nobreak\vspace{0.05em}\hspace*{1.5em}{\small #4}\par\vspace{0.42em}}
\begin{document}
\pagestyle{plain}

\begin{center}
{\LARGE\bfseries 动词时态语态复习 \quad 二刷 · 参考答案与解析}\\[0.25em]
{\large 2026 学年高一英语国庆专项（共 35 题）}\\[0.22em]
{\small 2029届高一英语 · \STAR 为 2026-10 原卷做错（10 题）及待核（第 17 题）}
\end{center}

\vspace{0.3em}
\noindent{\small 说明：题号沿用原卷。\STAR\ 为原卷做错的题（10 道），解析以"时态／语态信号词"为主线。}

\vspace{0.45em}
\qa{1}{}{C}{stepped、sat、began 三个并列谓语，都用过去式。}
\qa{2}{\STAR}{A}{\emph{so far}（到目前为止）$\to$ 现在完成时 haven't heard。}
\qa{3}{}{C}{"到场时电影已开演十分钟"：begin 是短暂动词，不能接 for ten minutes，改用 be on 表状态 $\to$ had been on。}
\qa{4}{}{A}{if 条件句用一般现在时表将来（noises 与 keep 是被动，are not kept）；主句用 will have to。}
\qa{5}{}{D}{"上次见 Jane 时她正在摘棉花"：主句 saw 为过去，从句用过去进行 was picking。}
\qa{6}{\STAR}{B}{\emph{soon} 表示"在过去看将来" $\to$ 过去将来时 would leave。}
\qa{7}{\STAR}{B}{时态呼应："我以为（过去）我（更早之前）丢了" $\to$ 主句 thought + 从句 had lost。}
\qa{8}{}{D}{before 从句用一般过去 got，主句动作发生更早 $\to$ 过去完成 had they known。}
\qa{9}{\STAR}{C}{"我不在这里长期工作，只是临时帮忙" $\to$ 现在进行时 am just helping out 表临时状态。}
\qa{10}{}{D}{before 从句用一般现在时表将来；collection 与 complete 是被动关系 $\to$ is completed。}
\qa{11}{}{B}{"正当她读报时奶奶睡着了"：从句用过去进行 was reading，主句用一般过去 fell。}
\qa{12}{}{B}{\emph{several times} 表示"到目前为止的经历" $\to$ 现在完成时 have met。}
\qa{13}{}{B}{"他（当时）正盯着空中发呆" $\to$ 过去进行时 was just staring。}
\qa{14}{}{D}{问"你见到我的眼镜了吗"，关注现在的结果 $\to$ 现在完成时 Have you seen。}
\qa{15}{}{B}{\emph{at the time}（当时）$\to$ 过去进行时 was working。}
\qa{16}{}{D}{反问"你难道还没见过他吗"，用现在完成时 Haven't you met him yet。}
\qa{17}{\STAR}{C}{\emph{since liberation} $\to$ 现在完成时；take place 不及物、无被动 $\to$ have taken place。（原卷此题答案处涂改不清，见订正表旁注）}
\qa{18}{}{B}{unless 条件句用一般现在时表将来；measures 与 take 是被动 $\to$ are taken。}
\qa{19}{}{A}{by the time 引导的时间状语从句用一般现在时表将来 $\to$ is。}
\qa{20}{\STAR}{B}{修饰人的表情用 -ed 分词 surprised（surprising 修饰物）；suggested 此处＝"表明"，从句用陈述语气 hadn't expected。}
\qa{21}{}{C}{\emph{so far} $\to$ 现在完成时 has appeared。}
\qa{22}{}{D}{陈述"现在具备的能力／习惯"，用一般现在时 play（since the new year 只管后半句）。}
\qa{23}{}{A}{"技术正飞速变化"（进行中的变化）$\to$ 现在进行时 is changing。}
\qa{24}{\STAR}{D}{问句是过去进行时（What were you doing），答句需描述"当时正在做" $\to$ was starting to take a shower。}
\qa{25}{\STAR}{B}{\emph{in the past five years} $\to$ 现在完成时 has spread。}
\qa{26}{}{D}{Never 置句首要倒装；到"那一刻"为止从未如此开心 $\to$ 过去完成 had I felt。}
\qa{27}{}{B}{\emph{always} + 进行时表"老是……（不满）" $\to$ are always watching。}
\qa{28}{\STAR}{B}{钱"快用完了"（尚未用完，正在走向用尽）$\to$ 现在进行时 is running out；run out 作"用完"是不及物动词，无被动。}
\qa{29}{}{C}{用 do 的过去式表强调："你确实发火了" $\to$ did lose。}
\qa{30}{}{B}{\emph{recent science}（近期的科学）$\to$ 现在完成时 has shown。}
\qa{31}{}{C}{盒子在"警察注意力被吸引"之前就已被放置 $\to$ 过去完成时的被动 had been placed。}
\qa{32}{\STAR}{C}{商店"即将关门"（所以清仓半价），用现在进行时表将来 $\to$ is closing down。}
\qa{33}{\STAR}{B}{"到那时已一直等了两小时，并且还会继续等" $\to$ 过去完成进行时 had been waiting。}
\qa{34}{}{A}{祈使句 + or + 将来时的结果句 $\to$ will never reach。}
\qa{35}{}{A}{"我的（雨衣）正挂在门后"（眼前的状态）$\to$ 现在进行时 is hanging。}

\vspace{0.3em}
\noindent{\color{rulegray}\rule{\textwidth}{0.5pt}}
\vspace{0.3em}
\noindent\textbf{复核记录}\hfill\small 完成重做后填写

\vspace{0.3em}
\noindent\begin{tabular}{@{}p{0.31\textwidth}@{}p{0.34\textwidth}@{}p{0.35\textwidth}@{}}
重做得分：\underline{\hspace{1.2cm}}\,/\,35
& 仍错误的题号：\underline{\hspace{3.2cm}}
& 下次复习日期：\underline{\hspace{2.0cm}} \\
\end{tabular}

\vspace{0.3em}
\noindent 本轮易错点：\checkmarkbox\ 现在完成时\quad\checkmarkbox\ 现在进行时（临时/即将）\quad\checkmarkbox\ 过去将来/过去完成进行\quad\checkmarkbox\ 分词与语态

\vspace{0.3em}
\noindent 复习状态：\checkmarkbox\ 已掌握\quad\checkmarkbox\ 需要巩固\quad\checkmarkbox\ 需重点复习

\end{document}
